\documentclass[12pt]{article}
\usepackage[T1]{fontenc}
\usepackage[utf8]{inputenc}
\usepackage{lmodern}
\usepackage{textcomp}
\usepackage{enumitem}
\usepackage{geometry}
\usepackage{graphicx}
\usepackage{amsmath, amssymb}
\geometry{margin=1in}
\begin{document}
\begin{center}
Political Science - Graduate Level Assessment 21 \
Total Marks: 40 \
Answer all questions.
\end{center}

\section*{Multiple Choice (10 marks)}
\begin{enumerate}[label=\arabic*.]
\item When considering economic tools for security policy, which of the following is the odd one out?
\begin{enumerate}[label=(\alph*)]
    \item Withdrawal of economic trade rights with the domestic market.
    \item Export controls protecting technological advantage and further foreign policy objectives.
    \item Control of munitions and arms sales.
    \item Import restrictions to protect a domestic market from foreign goods.
\end{enumerate}
\item According to Realists, what accounts for the onset of the Cold War?
\begin{enumerate}[label=(\alph*)]
    \item Ideological differences
    \item A power vacuum
    \item The expansionist nature of the Soviet Union
    \item Both b and c
\end{enumerate}
\item What drives US foreign policy according to the Marxist perspective?
\begin{enumerate}[label=(\alph*)]
    \item Economic protectionism
    \item Class consciousness
    \item The search for new markets
    \item Superstructure
\end{enumerate}
\item How did the Cold War context shape US perceptions of the Third World?
\begin{enumerate}[label=(\alph*)]
    \item The US ignored the Third World
    \item Local developments were viewed through a geopolitical lens
    \item The US unreservedly supported decolonization
    \item None of the above
\end{enumerate}
\item Within American politics, the power to accord official recognition to other countries belongs to
\begin{enumerate}[label=(\alph*)]
    \item the Senate.
    \item the president.
    \item the Secretary of State.
    \item the chairman of the Joint Chiefs.
\end{enumerate}
\end{enumerate}

\section*{True / False (10 marks)}
\textit{Answer True or False}
\begin{enumerate}[label=\arabic*.]
\setcounter{enumi}{5}
\item True or False: The responsibility to remain sovereign is not considered an element of the responsibility to protect framework.
\item True or False: Malware exclusively refers to viruses or worms that infect and damage computer files.
\item True or False: Regime type is not a reliable predictor of which states will acquire nuclear weapons.
\item True or False: Cyber-crime encompasses any crime that involves computers and networks in its execution or impact.
\item True or False: It is widely believed that fewer than 7 states in the international system have developed nuclear weapons capabilities.
\end{enumerate}

\section*{Long Form Response (20 marks)}
\textit{Answer each question with a clear and complete explanation.}
\begin{enumerate}[label=\arabic*.]
\setcounter{enumi}{10}
\item Discuss Wendt's Lockean culture of anarchy, highlighting its characteristics, historical context, and implications for state behavior and international relations.
\item Discuss the key policies that define a grand strategy of Offshore Balancing, emphasizing their implications for international relations and state security responsibilities.
\end{enumerate}

\vfill
\noindent\textit{End of Paper}
\end{document}